\documentclass[11pt]{article}
\usepackage[a4paper,margin=1in]{geometry}
\usepackage[T1]{fontenc}
\usepackage{helvet}
\renewcommand{\familydefault}{\sfdefault}
\usepackage{xcolor}
\usepackage{booktabs}
\usepackage{array}
\usepackage{colortbl}
\usepackage[font=footnotesize,labelformat=empty]{caption}
\usepackage{titlesec}
\usepackage{fancyhdr}
\usepackage{enumitem}
\usepackage{amsmath}
\usepackage{lastpage}
\usepackage[hidelinks]{hyperref}

% ---- Bayesian Capital brand palette ----
\definecolor{bcnavy}{RGB}{11,31,58}
\definecolor{bcblue}{RGB}{32,74,135}
\definecolor{bcgold}{RGB}{176,141,87}
\definecolor{bcgrey}{RGB}{90,98,112}
\definecolor{bcrule}{RGB}{210,214,220}
\definecolor{bckeep}{RGB}{232,237,244}
\definecolor{bcact}{RGB}{247,241,228}

\titleformat{\section}{\color{bcnavy}\Large\bfseries}{\thesection}{0.6em}{}
\titleformat{\subsection}{\color{bcblue}\large\bfseries}{\thesubsection}{0.6em}{}
\titlespacing*{\section}{0pt}{1.4em}{0.5em}

\pagestyle{fancy}\fancyhf{}
\renewcommand{\headrulewidth}{0pt}
\renewcommand{\footrulewidth}{0.4pt}
\renewcommand{\footrule}{\hbox to\headwidth{\color{bcrule}\leaders\hrule height \footrulewidth\hfill}}
\fancyfoot[L]{\footnotesize\color{bcgrey}Bayesian Capital}
\fancyfoot[C]{\footnotesize\color{bcgrey}Confidential --- Internal Research}
\fancyfoot[R]{\footnotesize\color{bcgrey}Page \thepage\ of \pageref{LastPage}}

\newcommand{\kpi}[2]{\textbf{#1}\,{\footnotesize\color{bcgrey}#2}}
\setlength{\parindent}{0pt}\setlength{\parskip}{0.55em}

\begin{document}

% ---------------- Title block ----------------
\begin{titlepage}
\vspace*{1.2cm}
{\color{bcgold}\rule{\textwidth}{2pt}}\\[0.4cm]
{\color{bcnavy}\Huge\bfseries Bayes / OU Capital Split}\\[0.35cm]
{\color{bcblue}\LARGE Findings Memorandum}\\[0.5cm]
{\color{bcgold}\rule{\textwidth}{1pt}}\\[0.6cm]
{\large\color{bcgrey} Hybrid Bayesian / Ornstein--Uhlenbeck Volatility Harvester}\\[0.2cm]
{\large\color{bcgrey} What the OU sleeve protects, and the per-stock efficiency-optimal split}\\[1.2cm]
{\color{bcnavy}\large\bfseries Bayesian Capital}\\[0.15cm]
{\color{bcgrey} Quantitative Research}\\[0.15cm]
{\color{bcgrey} 30 July 2026}\\[2cm]

\fbox{\parbox{0.92\textwidth}{\color{bcnavy}\textbf{One-line conclusion.}\\ \color{black}
The OU sleeve is a \emph{decorrelated volatility and drawdown dampener}, not an inverse hedge
and not crash protection. By the efficiency (Sharpe) metric the optimal split is
\emph{OU-heavy}, not Bayes-heavy --- but largely because OU sits in cash $\sim$91\% of the
time. The durable, regime-robust actions are per-name: \textbf{RKLB $\to$ OU-dominant} and
\textbf{AVGO $\to$ Bayes-heavy}, each keeping a 5\% floor on the minor sleeve so the full trade
count is preserved; most other names lean OU.}}
\end{titlepage}

% ---------------- Metrics primer ----------------
\section{How to read these numbers}

This memo judges each choice on three measures rather than on return alone. Return by itself
is misleading: a strategy can post a large number simply by taking large risk, and two books
with the same return can be worlds apart in how punishing they are to hold. The three
measures below separate \emph{how much} you earn from \emph{how much risk and pain} it costs.

\textbf{Sharpe ratio --- return per unit of volatility.} The Sharpe ratio divides a strategy's
average return by the variability (standard deviation) of those returns, scaled to an annual
figure. It answers: \emph{for every unit of ``bumpiness'' in the equity curve, how much return
am I paid?} A higher Sharpe means a smoother ride for the same return, or more return for the
same wobble. It is the standard yardstick for comparing strategies on a like-for-like,
risk-adjusted basis --- here, comparing the Bayes and OU sleeves and each capital split. Its
one blind spot: it penalises \emph{all} variability, including large \emph{upside} moves, and it
assumes risk is well described by day-to-day volatility.

\textbf{Maximum drawdown --- the worst peak-to-trough fall.} The maximum drawdown is the
largest drop from a previous high-water mark to a subsequent low, measured on the equity
curve. It is the loss you would actually have lived through at the worst moment --- the number
that determines whether you can stay invested, meet margin, and avoid capitulating at the
bottom. Where Sharpe treats up-moves and down-moves symmetrically, drawdown captures
\emph{downside path risk} directly: the depth of the hole. Two strategies can share a Sharpe
yet have very different drawdowns, which is why we always report it alongside.

\textbf{Calmar ratio --- return per unit of worst-case loss.} The Calmar ratio divides the
annualised return by the maximum drawdown. It answers: \emph{how much return do I earn for
each unit of worst-case pain?} Unlike Sharpe, it ignores harmless upside volatility and
penalises only the depth of the drawdown, so it suits a return-seeking book that cares about
surviving the worst loss rather than smoothing every daily wiggle. A Calmar of 3, for example,
means the strategy earned three times its deepest drawdown in annual return.

\textbf{Why all three.} Each metric is incomplete on its own --- return ignores risk, Sharpe
over-penalises good volatility, drawdown ignores return. Read together they triangulate the
real question: \emph{is the extra return worth the extra pain?} Optimising to any single number
invites a distorted answer, which is exactly the trap this memo is trying to avoid.

% ---------------- Executive summary ----------------
\section{Executive summary}

We asked two linked questions: \emph{when does the OU sleeve actually protect the book}, and
\emph{given Bayes earns far more, is there a case to hold most capital in Bayes?} All figures
are from the Python replica of the ten-name workbook, validated to reproduce every Model
sheet's cached results to the dollar.

\begin{itemize}[leftmargin=1.4em,itemsep=2pt]
  \item \kpi{OU is a dampener, not a hedge.}{} The sleeves are only $\sim$0.24 daily-correlated,
        and OU carries roughly \emph{half} the standalone drawdown of Bayes. It smooths the
        equity curve in normal and choppy conditions. It does \emph{not} rise when Bayes falls
        (it is up on only $\sim$24\% of down-Bayes days), and it co-moves with Bayes in a
        correlated sector crash --- the one time you might want protection.
  \item \kpi{Tilting to Bayes is a return-for-risk trade.}{} Book return rises monotonically
        with the Bayes weight ($344\%\!\to\!394\%$), but Sharpe collapses ($9.7\!\to\!2.6$) and
        max drawdown doubles ($5\%\!\to\!42\%$). There is no free lunch in the split.
  \item \kpi{Efficiency wants OU, with a catch.}{} The Sharpe-optimal split is $\sim$10--20\%
        Bayes for most names --- \emph{opposite} to a Bayes tilt. But OU's high Sharpe reflects
        that it is deployed only $\sim$9\% of days (Bayes $\sim$34\%); at an OU-heavy weight most
        capital simply sits in cash earning interest.
  \item \kpi{The two firm per-name calls.}{} \textbf{RKLB} should be OU-dominant (OU: 569\% at a
        9\% drawdown vs Bayes: 604\% at a 50\%) and \textbf{AVGO} Bayes-heavy --- the one name whose
        OU sleeve draws down \emph{more} than its Bayes sleeve. Keeping a 5\% floor on the minor
        sleeve in each preserves their full trade count at essentially no cost (\S6).
  \item \kpi{Trade count barely constrains the choice.}{} The number of buys is \emph{flat across
        every interior split} --- the split changes trade \emph{size}, not \emph{count} --- and
        falls only if a sleeve is zeroed. So the split can be set on Sharpe/Calmar alone, provided
        both sleeves stay funded.
\end{itemize}

% ---------------- Standalone ----------------
\section{When does the OU sleeve protect us?}

Run each sleeve standalone on the full name capital (frozen parameters):

\begin{table}[h]\centering\small
\renewcommand{\arraystretch}{1.12}
\begin{tabular}{l rrr c rrr c c}
\toprule
& \multicolumn{3}{c}{\textbf{Bayes}} & & \multicolumn{3}{c}{\textbf{OU}} & \textbf{daily} & \textbf{OU up on}\\
\cmidrule(lr){2-4}\cmidrule(lr){6-8}
\textbf{Name} & ann & Sh & maxDD & & ann & Sh & maxDD & \textbf{corr} & \textbf{Bayes-down day}\\
\midrule
NVDA & 203\% & 3.64 & 20.5\% & & 78\%  & 3.81 & 7.6\%  & 0.08 & 21\%\\
TSM  & 186\% & 3.93 & 21.1\% & & 59\%  & 5.65 & 3.4\%  & 0.21 & 17\%\\
TSLA & 163\% & 2.13 & 26.5\% & & 122\% & 2.97 & 18.1\% & 0.25 & 27\%\\
VRT  & 279\% & 3.52 & 29.1\% & & 96\%  & 3.66 & 13.6\% & 0.31 & 11\%\\
VST  & 312\% & 4.00 & 22.9\% & & 190\% & 5.40 & 11.4\% & 0.15 & 32\%\\
AVGO & 152\% & 3.68 & 14.5\% & & 132\% & 2.96 & 21.2\% & 0.15 & 14\%\\
PLTR & 189\% & 2.07 & 39.5\% & & 159\% & 3.79 & 20.3\% & 0.33 & 33\%\\
RKLB & 604\% & 2.46 & 49.9\% & & 569\% & 7.90 & 9.3\%  & 0.16 & 37\%\\
SOFI & 275\% & 3.51 & 22.8\% & & 179\% & 5.38 & 9.0\%  & 0.42 & 17\%\\
SPOT & 176\% & 4.70 & 9.4\%  & & 118\% & 7.79 & 5.0\%  & 0.32 & 31\%\\
\bottomrule
\end{tabular}
\caption*{\footnotesize Bayes wins raw return everywhere; OU wins Sharpe in 8/10 and roughly
halves drawdown. Correlation is low (0.08--0.42), but OU rises on only 11--37\% of the days
Bayes falls --- it diversifies, it does not counter.}
\end{table}

\textbf{The answer.} OU protects through \emph{low correlation and low volatility}: it holds
the book's Sharpe up and its drawdown down day-to-day, and it earns respectable returns at a
fraction of the risk. It does \emph{not} protect against a correlated event --- in the
December~2025 AI-sector drawdown the OU and Bayes sleeves stopped out together (see the
stop-loss study). OU is insurance against ordinary volatility and against the sideways/bear
regime this bull-market sample never tested --- not against a factor crash.

% ---------------- Split sweep ----------------
\section{Is there a case to hold most capital in Bayes?}

Sweeping the book-wide Bayes weight (allocation-weighted over the included names) lays the
trade-off bare:

\begin{table}[h]\centering\small
\renewcommand{\arraystretch}{1.12}
\begin{tabular}{c rrr l}
\toprule
\textbf{Bayes \%} & \textbf{book ann} & \textbf{Sharpe} & \textbf{maxDD} & \\
\midrule
0\%   & 344\% & 9.66 & 5.2\%  & all OU\\
25\%  & 357\% & 5.49 & 17.3\% & \\
50\%  & 370\% & 3.68 & 28.4\% & Model sheets\\
\rowcolor{bckeep} 60\%  & 375\% & 3.33 & 31.8\% & current book (flat mode)\\
70\%  & 380\% & 3.07 & 34.8\% & \\
80\%  & 385\% & 2.87 & 37.4\% & \\
100\% & 394\% & 2.59 & 41.7\% & all Bayes\\
\bottomrule
\end{tabular}
\caption*{\footnotesize Every step toward Bayes buys a little more return and a lot more risk.
100\% Bayes adds $\sim$19pp of return over the current 60\% but lifts max drawdown to 42\% and
delivers the book's worst Sharpe. A Bayes tilt is defensible only as a deliberate
return-seeking choice, and $\sim$65--70\% is the ceiling we would suggest.}
\end{table}

% ---------------- Per-stock optimal ----------------
\section{Per-stock efficiency-optimal split}

Because the sleeves are only $\sim$0.24 correlated, a \emph{blend} can beat either pure sleeve
on risk-adjusted terms. Gridding the split 0--100\% per name gives an interior optimum. We
report the Sharpe-max blend (pure efficiency) and the Calmar-max blend (return $\div$ drawdown,
which tolerates more Bayes), alongside each sleeve's capital deployment.

\begin{table}[h]\centering\small
\renewcommand{\arraystretch}{1.15}
\begin{tabular}{l cc c cc}
\toprule
& \multicolumn{2}{c}{\textbf{\% days deployed}} & & \textbf{Sharpe-opt} & \textbf{Calmar-opt}\\
\cmidrule(lr){2-3}\cmidrule(lr){5-6}
\textbf{Name} & Bayes & OU & & \textbf{Bayes \%} & \textbf{Bayes \%}\\
\midrule
NVDA & 28\% & 8\%  & & 20\% & 20\%\\
TSM  & 40\% & 1\%  & & 10\% & 10\%\\
TSLA & 49\% & 18\% & & 15\% & 25\%\\
VRT  & 27\% & 4\%  & & 20\% & 45\%\\
VST  & 26\% & 8\%  & & 15\% & 15\%\\
\rowcolor{bcact} AVGO & 21\% & 22\% & & 55\% & 100\%\\
PLTR & 55\% & 15\% & & 5\%  & 25\%\\
\rowcolor{bcact} RKLB & 43\% & 6\%  & & 0\%  & 0\%\\
SOFI & 35\% & 10\% & & 10\% & 5\%\\
SPOT & 15\% & 2\%  & & 10\% & 55\%\\
\bottomrule
\end{tabular}
\caption*{\footnotesize Both criteria are OU-leaning for almost every name. Gold rows are the
two exceptions: RKLB (OU-dominant) and AVGO (Bayes-heavy). These raw optima sit at the corners
(0\%/100\%); to preserve trade count we floor the minor sleeve at 5\% (\S6). Bayes\,\% is the
fraction of the name's capital in the Bayes tranche; OU\,\% is the remainder.}
\end{table}

\textbf{The catch.} OU is deployed only $\sim$9\% of days on average (TSM: 1\%), versus Bayes at
$\sim$34\%. Its high Sharpe is substantially a low-volatility, mostly-in-cash artifact. So
``Sharpe-optimal $=$ 80\% OU'' means parking most capital in a sleeve that trades a handful of
times a year. Efficiency and capital-utilisation pull in opposite directions; the split is a
choice of objective, not a single right number.

\subsection{Book-level effect}

\begin{table}[h]\centering\small
\renewcommand{\arraystretch}{1.12}
\begin{tabular}{l rrr}
\toprule
\textbf{Split policy} & \textbf{book ann} & \textbf{Sharpe} & \textbf{maxDD}\\
\midrule
Flat 50\% (Model sheets) & 370\% & 3.68 & 28.4\%\\
Flat 60\% (current book) & 375\% & 3.33 & 31.8\%\\
Per-stock Sharpe-opt     & 355\% & 6.23 & 14.4\%\\
Per-stock Calmar-opt     & 351\% & 9.83 & 4.6\%\\
\bottomrule
\end{tabular}
\caption*{\footnotesize Per-stock tilting gives up only $\sim$5--6\% of return versus the flat
60\% while roughly halving drawdown. \emph{The absolute Sharpe/drawdown magnitudes here are
in-sample, bull-regime artifacts and are not forward-achievable --- trust the ranking across
policies, not the levels.}}
\end{table}

% ---------------- Trade count ----------------
\section{Trade count and the split}

A track record is only as trustworthy as the number of independent trades behind it, so the
split must respect the buy count. The result is convenient: \textbf{the number of buys is flat
across every interior split and falls only at the corners.}

The reason is mechanical. A funded tranche places its bids and fills on price triggers that do
not depend on how much capital it holds; the split changes each trade's \emph{size}, not
\emph{whether} it happens. A 10/90 split therefore trades exactly as often as a 50/50 --- only
\emph{zeroing} a sleeve removes its trades:

\begin{table}[h]\centering\small
\renewcommand{\arraystretch}{1.1}
\begin{tabular}{l ccccc}
\toprule
\textbf{Bayes \%} & 0\% & 10\% & 50\% & 90\% & 100\%\\
\midrule
NVDA total buys & 47 & 184 & 184 & 184 & 137\\
\bottomrule
\end{tabular}
\caption*{\footnotesize Representative: identical buy count at 10/50/90\%; only the all-OU and
all-Bayes corners drop trades (losing the Bayes or the OU buys respectively).}
\end{table}

Two consequences. First, the split may be set purely on Sharpe/Calmar without sacrificing trade
count, as long as both sleeves stay funded. Second, the only recommendations that cost trades
are the two corner calls --- and both are rescued almost for free by a small floor:

\begin{table}[h]\centering\small
\renewcommand{\arraystretch}{1.15}
\begin{tabular}{l cc c cc}
\toprule
& \multicolumn{2}{c}{\textbf{corner}} & & \multicolumn{2}{c}{\textbf{5\% floor}}\\
\cmidrule(lr){2-3}\cmidrule(lr){5-6}
\textbf{Name} & split & buys / Sh / DD & & split & buys / Sh / DD\\
\midrule
RKLB & 0\% Bayes   & 183 / 7.90 / 9.3\%  & & 5\% Bayes  & \textbf{371} / 7.85 / 9.7\%\\
AVGO & 100\% Bayes & 142 / 3.68 / 14.5\% & & 90\% Bayes & \textbf{263} / 3.86 / 15.0\%\\
\bottomrule
\end{tabular}
\caption*{\footnotesize A 5\% floor on the minor sleeve restores the full trade count at
negligible cost --- RKLB's Sharpe and drawdown barely move; AVGO's Sharpe actually improves.}
\end{table}

\textbf{Rule: never run a corner --- floor both sleeves at 5\%.} This keeps every name's full
buy count while leaving the interior Sharpe/Calmar choice untouched.

% ---------------- Implementation ----------------
\section{Recommendation and implementation}

Populate the Allocation sheet's \textbf{split mode~2} (per-stock Bayes\,\%) with one of the two
columns above, by objective:

\begin{itemize}[leftmargin=1.4em,itemsep=2pt]
  \item \kpi{Efficiency first (recommended):}{} the \textbf{Sharpe-opt} column, floored ---
        NVDA~20, TSM~10, TSLA~15, VRT~20, VST~15, AVGO~55, PLTR~5, RKLB~\textbf{5}, SOFI~10,
        SPOT~10.
  \item \kpi{Return-aware:}{} the \textbf{Calmar-opt} column, floored, if you want capital
        working harder and can accept more drawdown --- differences are VRT~45, TSLA~25, PLTR~25,
        AVGO~\textbf{90}, SPOT~55, RKLB~\textbf{5}; the rest unchanged.
\end{itemize}

Both columns apply the 5\% floor of \S6 (RKLB~5 not 0; AVGO~90 not 100), so the full trade count
is retained. Make the two firm calls first: \textbf{RKLB $\to$ OU-dominant (5\% Bayes)} and
\textbf{AVGO $\to$ Bayes-heavy (90\%)}. They are the largest, most robust corrections to the
uniform 50/50 and hold under both criteria.

\textbf{Caveats.} (1) All figures are in-sample over a strong bull regime for these names; the
book-level Sharpe/drawdown levels are inflated and only the cross-name ranking is trustworthy.
(2) The split, like the stop period, is a single low-dimensional knob --- use the round numbers
above, not razor-edge argmaxima. (3) A Bayes tilt is a bet that the trending regime persists; in
a sideways or bear market the OU-lean would look better still.

\vfill
{\footnotesize\color{bcgrey} Data vintage: current ten-name workbook (\texttt{\_nopq} PLTR
build), Apr 2024 onward, 2.2-year horizon. Engine reproduces every Model sheet's cached results
to the dollar. Split values are fractions of each name's own capital; the remainder goes to OU.}

\end{document}
